\documentclass[10pt,a4paper]{amsart}
\usepackage[margin=19mm]{geometry}
\usepackage{amsmath,amssymb,mathtools}
\usepackage[scr=rsfs,cal=euler]{mathalfa}
\usepackage{xfrac}
\usepackage[unicode,hidelinks,bookmarks=false]{hyperref}
\newcommand{\ie}{i.e.\ }
\newcommand{\cf}{cf.\ }
\newcommand{\resp}{respectively\ }
\newcommand{\Subsection}{\subsection}
\providecommand{\nobreakdash}{\nobreak\hbox{-}}
\setlength{\emergencystretch}{2em}
\multlinegap0pt
\usepackage{amssymb}
\usepackage{bm}
\usepackage{booktabs}
\usepackage{xcolor}
\newtheorem{lemma}{Lemma}
\title{Quadratic Embedding Constants of Corona Graphs}
\author{Ferdi, Edy Tri Baskoro, Nobuaki Obata, Aditya Purwa Santika}
\date{Source: arXiv:2512.17258v1 (CC BY 4.0)}
\begin{document}
\maketitle

\noindent\emph{Source and licence.} Quadratic Embedding Constants of Corona Graphs. Version arXiv:2512.17258v1 (CC BY 4.0), distributed by arXiv under the Creative Commons Attribution 4.0 International licence, http://creativecommons.org/licenses/by/4.0/, as recorded in the arXiv licence metadata for this exact version and shown on the abstract page https://arxiv.org/abs/2512.17258v1. Full licence notice: \texttt{licenses/arxiv-2512.17258-CC-BY-4.0.txt}.\par\smallskip

\noindent\textit{Editorial scope.} Counted unit: the article's lemma (1) (source0.tex lines 981--994). The article's complete original proof of it is delivered with it (source0.tex lines 996--1071, one \texttt{proof} environment of 76 lines). No mathematics is written, repaired or re-derived: the statement and the proof are the article's own lines, in the article's own notation, and no retained line is substituted or re-flowed.  Retained uncounted as the source-local context the proof needs: lines 922--926; lines 932--959.  Nothing was omitted from any retained span, so the package carries no omission record.  No retained line contains a citation, so the wrapper carries no bibliography and no bibliography entry.  No retained line contains a figure environment, an image-inclusion command or drawing code, so no image is delivered and the package declares no asset dependency.  Editorial formatting only, in addition: an AMS \texttt{amsart} preamble that restates the article's own declarations and macros used by the retained lines, namely the environments \texttt{lemma} on separate counters, together with the standard packages those lines need.  The retained environments print in this document as Lemma 1 in order of appearance, whereas the article prints the same items with the numbering of its own full document.  The complete native source of the article as distributed in the arXiv e-print for this exact version is bundled unchanged as \texttt{sources/191304/source0.tex}.
\begin{equation}\label{03eqn:omega(lambda) as rational function}
\omega_A(\lambda)
=p(\lambda)\prod_{\alpha\in\mathrm{ev}_{\mathrm{M}}(A)\backslash\{-2\}}
 (\alpha+2+\lambda)^{-1},
\end{equation}

\begin{lemma}\label{03lem:behaviour of omega(lambda)}
Let $A\in M_n(\mathbb{R})$ be a symmetric matrix
and $\omega_A(\lambda)$ the associated $\omega$-function
defined as above. 
%
\begin{enumerate}
\item[\upshape (1)] 
$\omega_A(0)=1+\|E_{-2}\mathbf{1}\|^2$.
\item[\upshape (2)] 
$\displaystyle\lim_{\lambda\rightarrow\pm\infty}\omega_A(\lambda)=n+1$.
\item[\upshape (3)]
For $\alpha\in\mathrm{ev}_{\mathrm{M}}(A)\backslash\{-2\}$ with $\alpha+2>0$
we have
\[
\lim_{\lambda\downarrow -\alpha-2}\omega_A(\lambda)=-\infty,
\qquad
\lim_{\lambda\uparrow -\alpha-2}\omega_A(\lambda)=+\infty.
\]
\item[\upshape (4)]
For $\alpha\in\mathrm{ev}_{\mathrm{M}}(A)\backslash\{-2\}$ with $\alpha+2<0$
we have
\[
\lim_{\lambda\downarrow -\alpha-2}\omega_A(\lambda)=+\infty,
\qquad
\lim_{\lambda\uparrow -\alpha-2}\omega_A(\lambda)=-\infty.
\]
\end{enumerate}
\end{lemma}

\begin{lemma}\label{03lem:zeroes of omega(lambda) (1)}
Assume that $k_A\ge1$ and $-2\not\in\mathrm{ev}_{\mathrm{M}}(A)$.
Arrange the members in $\mathrm{ev}_{\mathrm{M}}(A)\cup \{-2\}$
in ascending order as
\[
\mathrm{ev}_{\mathrm{M}}(A)\cup \{-2\}
=\{\alpha_1<\alpha_2<\dots<\alpha_{k_A+1}\}.
\]
Then, for each $1\le i\le k_A$ the interval
$(-\alpha_{i+1}-2,-\alpha_i-2)$
contains a unique zero of $\omega_A(\lambda)$.
Moreover, those $k_A$ distinct real zeroes
exhaust the zeroes of $\omega_A(\lambda)$.
\end{lemma}

\begin{proof}
For simplicity we write $k=k_A$.
Since $\omega_A(\lambda)$ is a rational function
as in \eqref{03eqn:omega(lambda) as rational function}
with $p(\lambda)$ being a polynomial of order $k$,
$\omega_A(\lambda)$ has at most $k$ distinct zeroes
in the complex plane.
Therefore, for the assertion it is sufficient to show that
each interval $(-\alpha_{i+1}-2,-\alpha_i-2)$
for $1\le i\le k$
contains at least one zero of $\omega_A(\lambda)$.

(Case 1) $\alpha_1=-2$. 
In this case, we have $-\alpha_k-2<\dotsb<-\alpha_2-2<0$.
By definition,
$\omega_A(\lambda)$ is continuous on the interval
$(-\alpha_{i+1}-2,-\alpha_i-2)$ for each $2\le i\le k$.
Moreover, 
we see from Lemma \ref{03lem:behaviour of omega(lambda)} (3) that
\[
\lim_{\lambda\downarrow -\alpha_{i+1}-2}\omega_A(\lambda)=-\infty,
\qquad
\lim_{\lambda\uparrow -\alpha_i-2}\omega_A(\lambda)=+\infty.
\]
Therefore, $\omega_A(\lambda)$ has at least one zero
in the interval $(-\alpha_{i+1}-2,-\alpha_i-2)$ for each $2\le i\le k$.
Moreover, since 
$\omega_A(\lambda)$ is a continuous function on
$(-\alpha_2-2,+\infty)$ and
\[
\lim_{\lambda\downarrow -\alpha_2-2}\omega(\lambda)=-\infty,
\qquad
\omega_A(0)=1,
\]
it has at least one zero
in the interval $(-\alpha_2-2,0)=(-\alpha_2-2,-\alpha_1-2)$.

(Case 2) $\alpha_k=-2$. In a similar manner as in (Case 1),
we may show that
for each $1\le i\le k$ the interval $(-\alpha_{i+1}-2,-\alpha_i-2)$
contains at least one zero of $\omega_A(\lambda)$.


(Case 3) $\alpha_1<-2<\alpha_k$. In that case we have
$\alpha_l=-2$ for some $2\le l \le k-1$.
For $1\le i\le l-2$ the function $\omega_A(\lambda)$ is continuous
on $(-\alpha_{i+1}-2,-\alpha_i-2)$ and
\[
\lim_{\lambda\downarrow -\alpha_{i+1}-2}\omega_A(\lambda)=+\infty,
\qquad
\lim_{\lambda\uparrow -\alpha_i-2}\omega_A(\lambda)=-\infty.
\]
by Lemma \ref{03lem:behaviour of omega(lambda)} (4).
Therefore, $\omega_A(\lambda)$ has at least one zero
in the $(-\alpha_{i+1}-2,-\alpha_i-2)$ for $1\le i\le l-2$.
Similarly, $\omega_A(\lambda)$ has at least one zero
in the $(-\alpha_{i+1}-2,-\alpha_i-2)$ for
$l+1\le i\le k$.
Finally, since $\omega_A(\lambda)$ is continuous
on $(-\alpha_{l+1}-2,-\alpha_{l-1}-2)$ and
\[
\lim_{\lambda\downarrow -\alpha_{l+1}-2}\omega_A(\lambda)=-\infty,
\qquad
\omega_A(0)=1,
\qquad
\lim_{\lambda\uparrow -\alpha_{l-1}-2}\omega_A(\lambda)=-\infty
\]
by Lemma \ref{03lem:behaviour of omega(lambda)},
it has at least one zero
in each of the intervals
$(-\alpha_{l+1}-2,0)=(-\alpha_{l+1}-2,-\alpha_l-2)$
and $(0,-\alpha_{l-1}-2)=(-\alpha_l-2,-\alpha_{l-1}-2)$.
Consequently, $\omega_A(\lambda)$ 
has at least one zero in each interval $(-\alpha_{i+1}-2,-\alpha_i-2)$
for $1\le i\le k$.
\end{proof}

\end{document}
